\documentclass[../../main]{subfiles}
\begin{document}

This chapter introduces the theoretical background required for the remainder of the thesis.

The first section presents the fundamentals of Constraint Programming (CP), including CP solvers and the MiniZinc modeling language used in the experimental evaluation.  
The second section introduces Large Language Models (LLMs), focusing on their probabilistic formulation, inference mechanisms, and prompting strategies.

\section{Constraint Programming}

Constraint Programming (CP) is a paradigm for declaratively modeling and solving combinatorial problems, particularly in domains such as planning and scheduling~\cite{CPholygrail}.

A constraint is a logical relation defined over a set of variables, each associated with a domain of possible values. Constraints restrict the combinations of values that variables may simultaneously assume and therefore represent partial information about the system being modeled. An important feature of constraints is their declarative nature: they specify relationships that must hold without prescribing the procedure used to enforce them~\cite{CPholygrail}.

CP studies computational systems that solve problems by enforcing sets of such constraints and searching for assignments of variable values that satisfy them.

\subsection{Solvers}

In this thesis, a solver denotes a system that computes assignments satisfying a set of constraints and, when required, optimizes an objective function.

A Constraint Satisfaction Problem (CSP) can be formalized as the tuple
\[
\langle X, D, C \rangle,
\]
where \(X\) is a finite set of decision variables, \(D(x)\) is the domain of each variable \(x \in X\), and \(C\) is a set of constraints defined over subsets of \(X\). A solution is an assignment
\[
s : X \rightarrow \bigcup_{x \in X} D(x)
\]
such that \(s(x) \in D(x)\) for all \(x \in X\) and all constraints in \(C\) are satisfied~\cite{SeaPearl}.

A Constraint Optimization Problem (COP) extends the CSP by introducing an objective function. Formally, a COP is defined as
\[
\langle X, D, C, f \rangle,
\]
where \(f : S \rightarrow \mathbb{R}\) is defined over the set \(S\) of feasible assignments. The objective is to compute an optimal solution
\[
s^* = \arg\min_{s \in S} f(s)
\quad \text{or} \quad
s^* = \arg\max_{s \in S} f(s).
\]

Different solver families address CSPs and COPs using distinct mathematical abstractions and algorithmic techniques.

\subsubsection{Constraint Programming solvers}

Constraint Programming solvers operate directly on variables with finite or structured domains and heterogeneous constraints. Their core mechanisms are constraint propagation, which removes inconsistent values from variable domains, and systematic search, typically implemented through backtracking guided by heuristics. CP solvers also support high-level global constraints equipped with specialized filtering algorithms that exploit combinatorial structure~\cite{CPvsMILPsolvers}. A representative example is Gecode~\cite{Gecode}.

\subsubsection{Portfolio solvers}

Solving combinatorial search problems is computationally challenging, and different solvers often perform better on different instances of the same problem class. Algorithm portfolios exploit this observation by combining multiple solvers and selecting the most appropriate one for each instance~\cite{PortfolioSolvers}.

Algorithm portfolios are a specific case of the broader Algorithm Selection problem~\cite{AlgSelProb}. In a portfolio solver, a set of constituent solvers \(s_1, \dots, s_m\) is available. Given a problem instance \(p\), the system predicts which solver (or subset of solvers) is most suitable for solving it.

Solver selection is typically performed using machine learning models trained on features extracted from the problem instance~\cite{AlgSelCSP}. In this work, we investigate whether a pre-trained LLM can serve as the decision component for this selection process, potentially avoiding the need for dedicated training~\cite{fromPFStoAS}.

\subsection{MiniZinc}

MiniZinc is a high-level modeling language for constraint programming that supports a wide range of solvers and enables solver-independent problem specification~\cite{MiniZinc}. It was designed to provide a standard modeling language for CP and to facilitate solver benchmarking.

For these reasons, all problems used in the experimental evaluation were modeled in MiniZinc.

\subsubsection{A MiniZinc example}

A MiniZinc specification consists of a model, describing the structure of a class of problems, and a data component specifying a particular instance. Data files contain assignments to parameters declared in the model and are typically provided separately.

A MiniZinc model is composed of a sequence of items, including variable declarations, constraints, predicates, and a single \texttt{solve} item specifying whether the problem is a satisfaction or optimization task.

\Cref{fig:minizincModelEx} and \Cref{fig:minizincDataEx} show an example MiniZinc model and data file for a job shop scheduling problem~\cite{MiniZinc}. The model defines parameters, decision variables, and constraints that enforce task ordering and prevent resource conflicts, while the data file provides the specific instance parameters.

\begin{figure}[ht]
    \centering
    \includegraphics[width=0.75\linewidth]{minizincModelEx.png}
    \caption{MiniZinc model (\texttt{jobshop.mzn}) for the job shop problem (adapted from~\cite{MiniZinc}).}
    \label{fig:minizincModelEx}
\end{figure}

\begin{figure}[ht]
    \centering
    \includegraphics[width=0.75\linewidth]{minizincDataEx.png}
    \caption{MiniZinc data (\texttt{jobshop2x2.dzn}) for the job shop problem (adapted from~\cite{MiniZinc}).}
    \label{fig:minizincDataEx}
\end{figure}

\subsubsection{MiniZinc Challenge}

Comparing constraint programming systems is inherently difficult because different solvers support different modeling features and optimization strategies. MiniZinc addresses this issue by providing a standardized modeling language that allows multiple solvers to be evaluated on the same problem specifications.

Since 2008, the MiniZinc Challenge has provided a yearly benchmark competition comparing solvers that support MiniZinc~\cite{PhilMznChallenge}. Participants submit solvers capable of handling MiniZinc or FlatZinc models, which are then evaluated on a standardized set of benchmark instances. Performance is assessed based on solution quality, number of solved instances, and solving time.

Although the MiniZinc Challenge provides a comparative scoring system, this thesis employs alternative evaluation metrics better suited to absolute performance assessment in solver-selection experiments.

\section{Large Language Models}
\label{sec:LLM}

Large Language Models (LLMs) are statistical language models that estimate probability distributions over sequences of tokens by learning patterns from large text corpora.

Modern LLMs are typically implemented using transformer architectures~\cite{TransformerArchitecture}, which rely on self-attention mechanisms to model dependencies between tokens in a sequence. Given an input sequence \(x_1, \dots, x_n\), the model estimates the conditional probability

\[
P(x_{n+1} \mid x_1, \dots, x_n).
\]

Text generation proceeds autoregressively by repeatedly sampling tokens from this distribution.

Training generally involves two stages. During large-scale pretraining, models learn general linguistic and semantic representations through self-supervised objectives. Alignment techniques, such as supervised fine-tuning and reinforcement learning from feedback, subsequently adapt the model to follow instructions and produce outputs consistent with user expectations~\cite{LLMsIBM}.

\subsection{Inference and decoding}

During inference, the model produces logits over the vocabulary, which are converted into probabilities using the softmax function~\cite{InfeerenceTrainingLLMs}. Decoding strategies determine how the next token is selected from this distribution.

Deterministic decoding methods include greedy search, which selects the most probable token at each step, and beam search~\cite{BeamSearch}, which maintains multiple candidate sequences during generation.

Stochastic methods introduce controlled randomness into the generation process. Temperature sampling modifies the probability distribution using a temperature parameter \(\tau\), while top-\(p\) sampling restricts token selection to the smallest subset of tokens whose cumulative probability exceeds a threshold \(p\)~\cite{LLMDecoding}.

A practical constraint is the context window, which limits the total number of tokens that can be processed simultaneously. This limitation influences prompt design and the amount of contextual information that can be provided to the model.

\subsection{Relevance to solver selection}

In this work, LLMs are used as decision components for solver selection. Given a representation of a problem instance, the model is prompted to recommend the most appropriate solver from a predefined portfolio.

The effectiveness of this approach depends both on the model's capacity to interpret structural properties of constraint models and on practical aspects such as prompt formulation, representation choice, and inference parameters.

\end{document}